% GENERATED FILE. Edit canonical lessons, then run npm run generate:books.
% canonical-source-hash: fnv1a64:6ff51b3a53556b3b
% canonical-lessons: TE-C375-patrikeyudu, TE-C375-mudrana, TE-C375-sancalanam, TE-C375-nivedincu, TE-C375-spandincu

\chapter{The Press}
\label{ch:te-the-press}
\hlchaptermodality{375}

\begin{chapteropening}
\textbf{By the end of this chapter:} \emph{I can say a journalist, printing, a sensation, to submit a report and to respond.}

Complete the last lesson of chapter 375: I can say a journalist, printing, a sensation, to submit a report and to respond.
\end{chapteropening}

\section[\texorpdfstring{\te{పాత్రికేయుడు} (pātrikēyuḍu)}{పాత్రికేయుడు (pātrikēyuḍu)}]{\te{పాత్రికేయుడు} (pātrikēyuḍu) — a journalist}

\label{lesson:TE-C375-patrikeyudu}



\begin{quote}
Before the new one: say the Telugu for to behold, then the Telugu for a face-to-face interview.
\end{quote}

\subsection*{You'll want to know: \te{పాత్రికేయుడు}}
\textbf{\te{పాత్రికేయుడు}} — \emph{pātrikēyuḍu} — \textquotedblleft{}a journalist\textquotedblright{}.

\begin{quote}
\te{ఆ} \te{పాత్రికేయుడు} \te{ముప్పై} \te{ఏళ్ళు} \te{పని} \te{చేశాడు}.
\end{quote}

\begin{quote}
\emph{ā pātrikēyuḍu muppai ēḷḷu pani cēśāḍu.}
\end{quote}

\begin{quote}
That journalist worked for thirty years.
\end{quote}

Plural: \te{పాత్రికేయులు}, \emph{pātrikēyulu}.

Built on \te{పత్రిక}, a periodical.

\begin{cousinweb}
\textbf{\te{పాత్రికేయుడు}} (\emph{pātrikēyuḍu}) is a modern word built on Sanskrit \emph{pattrikā}, \textquotedblleft{}a letter, a document\textquotedblright{} (Monier-Williams), the source of \te{పత్రిక}.
\end{cousinweb}

\begin{grammarlens}[title={What you've built}]
One more word for telling what happened.
\end{grammarlens}

\subsection*{Your turn}
\begin{itemize}
\raggedright
  \item \emph{Say it:} \emph{pātrikēyuḍu}
  \item \emph{Say it:} \emph{pātrikēyuḍu}, once more
  \item \emph{Say it:} \emph{ā pātrikēyuḍu muppai ēḷḷu pani cēśāḍu.} — That journalist worked for thirty years.
  \item \emph{Recall:} say the Telugu for to behold, then the Telugu for a face-to-face interview, then say \emph{pātrikēyuḍu} again
\end{itemize}

\begin{tcolorbox}[breakable,colback=teal!4,colframe=teal!35!black,title={Before you move on}]
What does \te{పాత్రికేయుడు} mean? (\textquotedblleft{}A journalist\textquotedblright{}.) Say it once more.
\end{tcolorbox}
\section[\texorpdfstring{\te{ముద్రణ} (mudraṇa)}{ముద్రణ (mudraṇa)}]{\te{ముద్రణ} (mudraṇa) — printing; an edition}

\label{lesson:TE-C375-mudrana}



\begin{quote}
Before the new one: say the Telugu for a face-to-face interview, then the Telugu for a journalist.
\end{quote}

\subsection*{You'll want to know: \te{ముద్రణ}}
\textbf{\te{ముద్రణ}} — \emph{mudraṇa} — \textquotedblleft{}printing; an edition\textquotedblright{}.

\begin{quote}
\te{ఆ} \te{పుస్తకం} \te{మూడో} \te{ముద్రణ} \te{వచ్చింది}.
\end{quote}

\begin{quote}
\emph{ā pustakaṁ mūḍō mudraṇa vaccindi.}
\end{quote}

\begin{quote}
The book's third edition came out.
\end{quote}

You know \te{ముద్రించు}, to print; \te{ముద్రణ} is the printing, or an edition.

\begin{cousinweb}
\textbf{\te{ముద్రణ}} (\emph{mudraṇa}) is borrowed from Sanskrit \emph{mudraṇa}, \textquotedblleft{}the act of sealing up\textquotedblright{} (CDIAL 10201), from \emph{mudrā}, \textquotedblleft{}a seal\textquotedblright{}.
\end{cousinweb}

\begin{grammarlens}[title={What you've built}]
One more word for telling what happened.
\end{grammarlens}

\subsection*{Your turn}
\begin{itemize}
\raggedright
  \item \emph{Say it:} \emph{mudraṇa}
  \item \emph{Say it:} \emph{mudraṇa}, once more
  \item \emph{Say it:} \emph{ā pustakaṁ mūḍō mudraṇa vaccindi.} — The book's third edition came out.
  \item \emph{Recall:} say the Telugu for a face-to-face interview, then the Telugu for a journalist, then say \emph{mudraṇa} again
\end{itemize}

\begin{tcolorbox}[breakable,colback=teal!4,colframe=teal!35!black,title={Before you move on}]
What does \te{ముద్రణ} mean? (\textquotedblleft{}Printing; an edition\textquotedblright{}.) Say it once more.
\end{tcolorbox}
\section[\texorpdfstring{\te{సంచలనం} (sañcalanaṁ)}{సంచలనం (sañcalanaṁ)}]{\te{సంచలనం} (sañcalanaṁ) — a sensation}

\label{lesson:TE-C375-sancalanam}



\begin{quote}
Before the new one: say the Telugu for a journalist, then the Telugu for printing.
\end{quote}

\subsection*{You'll want to know: \te{సంచలనం}}
\textbf{\te{సంచలనం}} — \emph{sañcalanaṁ} — \textquotedblleft{}a sensation\textquotedblright{}.

\begin{quote}
\te{ఆ} \te{వార్త} \te{రాష్ట్రమంతా} \te{సంచలనం} \te{సృష్టించింది}.
\end{quote}

\begin{quote}
\emph{ā vārta rāṣṭramantā sañcalanaṁ sṛṣṭiñcindi.}
\end{quote}

\begin{quote}
That news caused a sensation across the state.
\end{quote}

\te{సంచలనం} \te{సృష్టించు}, to cause a sensation.

\begin{cousinweb}
\textbf{\te{సంచలనం}} (\emph{sañcalanaṁ}) is borrowed from Sanskrit \emph{saṁcalana}, \textquotedblleft{}moving about, agitation, trembling\textquotedblright{} (Monier-Williams).
\end{cousinweb}

\begin{grammarlens}[title={What you've built}]
One more word for telling what happened.
\end{grammarlens}

\subsection*{Your turn}
\begin{itemize}
\raggedright
  \item \emph{Say it:} \emph{sañcalanaṁ}
  \item \emph{Say it:} \emph{sañcalanaṁ}, once more
  \item \emph{Say it:} \emph{ā vārta rāṣṭramantā sañcalanaṁ sṛṣṭiñcindi.} — That news caused a sensation across the state.
  \item \emph{Recall:} say the Telugu for a journalist, then the Telugu for printing, then say \emph{sañcalanaṁ} again
\end{itemize}

\begin{tcolorbox}[breakable,colback=teal!4,colframe=teal!35!black,title={Before you move on}]
What does \te{సంచలనం} mean? (\textquotedblleft{}A sensation\textquotedblright{}.) Say it once more.
\end{tcolorbox}
\section[\texorpdfstring{\te{నివేదించు} (nivēdiñcu)}{నివేదించు (nivēdiñcu)}]{\te{నివేదించు} (nivēdiñcu) — to submit a report}

\label{lesson:TE-C375-nivedincu}



\begin{quote}
Before the new one: say the Telugu for printing, then the Telugu for a sensation.
\end{quote}

\subsection*{You'll want to know: \te{నివేదించు}}
\textbf{\te{నివేదించు}} — \emph{nivēdiñcu} — \textquotedblleft{}to submit a report\textquotedblright{}.

\begin{quote}
\te{విలేకరి} \te{ఆ} \te{సంఘటన} \te{గురించి} \te{నివేదించాడు}.
\end{quote}

\begin{quote}
\emph{vilēkari ā saṅghaṭana guriñci nivēdiñcāḍu.}
\end{quote}

\begin{quote}
The reporter reported on the incident.
\end{quote}

Past (I …): \te{నివేదించాను}, \emph{nivēdiñcānu}.

You know the noun \te{నివేదిక}, a report.

\begin{cousinweb}
\textbf{\te{నివేదించు}} (\emph{nivēdiñcu}) is built on Sanskrit \emph{nivedana}, \textquotedblleft{}announcing, proclaiming\textquotedblright{} (Monier-Williams), with the Telugu verb ending \emph{-iñcu}.
\end{cousinweb}

\begin{grammarlens}[title={What you've built}]
One more word for telling what happened.
\end{grammarlens}

\subsection*{Your turn}
\begin{itemize}
\raggedright
  \item \emph{Say it:} \emph{nivēdiñcu}
  \item \emph{Say it:} \emph{nivēdiñcu}, once more
  \item \emph{Say it:} \emph{vilēkari ā saṅghaṭana guriñci nivēdiñcāḍu.} — The reporter reported on the incident.
  \item \emph{Recall:} say the Telugu for printing, then the Telugu for a sensation, then say \emph{nivēdiñcu} again
\end{itemize}

\begin{tcolorbox}[breakable,colback=teal!4,colframe=teal!35!black,title={Before you move on}]
What does \te{నివేదించు} mean? (\textquotedblleft{}To submit a report\textquotedblright{}.) Say it once more.
\end{tcolorbox}
\section[\texorpdfstring{\te{స్పందించు} (spandiñcu)}{స్పందించు (spandiñcu)}]{\te{స్పందించు} (spandiñcu) — to respond; to react}

\label{lesson:TE-C375-spandincu}



\begin{quote}
Before the new one: say the Telugu for a sensation, then the Telugu for to submit a report.
\end{quote}

\subsection*{You'll want to know: \te{స్పందించు}}
\textbf{\te{స్పందించు}} — \emph{spandiñcu} — \textquotedblleft{}to respond; to react\textquotedblright{}.

\begin{quote}
\te{పాఠకులు} \te{ఆ} \te{కథనానికి} \te{వెంటనే} \te{స్పందించారు}.
\end{quote}

\begin{quote}
\emph{pāṭhakulu ā kathanāniki veṇṭanē spandiñcāru.}
\end{quote}

\begin{quote}
Readers responded to the story at once.
\end{quote}

Past (I …): \te{స్పందించాను}, \emph{spandiñcānu}.

What you respond to takes -\te{కి}.

\begin{cousinweb}
\textbf{\te{స్పందించు}} (\emph{spandiñcu}) is built on the Sanskrit root \emph{spand}, \textquotedblleft{}to quiver, throb\textquotedblright{} (Monier-Williams), with the Telugu verb ending \emph{-iñcu}: to stir in answer.
\end{cousinweb}

\begin{grammarlens}[title={What you've built}]
One more word for telling what happened.
\end{grammarlens}

\subsection*{Your turn}
\begin{itemize}
\raggedright
  \item \emph{Say it:} \emph{spandiñcu}
  \item \emph{Say it:} \emph{spandiñcu}, once more
  \item \emph{Say it:} \emph{pāṭhakulu ā kathanāniki veṇṭanē spandiñcāru.} — Readers responded to the story at once.
  \item \emph{Recall:} say the Telugu for a sensation, then the Telugu for to submit a report, then say \emph{spandiñcu} again
\end{itemize}

\begin{tcolorbox}[breakable,colback=teal!4,colframe=teal!35!black,title={Before you move on}]
What does \te{స్పందించు} mean? (\textquotedblleft{}To respond; to react\textquotedblright{}.) Say it once more.
\end{tcolorbox}
